\begin{proposition}[Proved: a generically finite morphism to a projective target]
\label{D:generically-finite-projective-target}
Let $X$ be a normal connected proper scheme over an
algebraically closed field.  Suppose $X$ admits a morphism
to a projective scheme that is generically finite onto its
image.  Then $Q(X,Y)$ is topologically finitely generated
for every proper closed subset $Y\subset X$.
In particular, if $Y$ has codimension at least two, the
full group $\pi_1(X\setminus Y)$ is topologically finitely
generated.
\end{proposition}
\begin{proof}
Replace the projective target by the reduced image.
Stein factorization gives a proper birational morphism
\[
                              q:X\longrightarrow Z,
\]
where $Z$ is normal and finite over that image, hence
projective.  More explicitly, the generic fibre of the
original map is the spectrum of the finite extension
$k(X)$ of the image function field.  Its Stein factor
is the normalization of the image in $k(X)$, so the
remaining morphism is birational.  We use
\cite[Theorem 37.53.4, tag 03H0]{C_StacksFibration}.

Let $Z^\circ$ be the largest open over which $q$ is an
isomorphism.  It contains all codimension-zero and
codimension-one points of the normal $Z$, by
\cite[Lemma 33.17.3, tag 0BFP]{C_StacksProjectiveTarget}.
Put $E=Z\setminus Z^\circ$ and
\[
                  W=Z\setminus\bigl(E\cup q(Y)\bigr).
\]
The map $q$ is proper, so $q(Y)$ is closed.  It is a proper
subset of $Z$, since $q$ is birational and $Y$ is proper
in the integral $X$.  Thus $W$ is a nonempty open, and
$q^{-1}(W)\simeq W$ is a dense open of $U=X\setminus Y$.

Every connected finite \'etale cover of the normal
integral $U$ remains connected on this dense open.
Moreover, a codimension-one component of
$Z\setminus W$ cannot be contained in $E$, since $E$
has codimension at least two.  Its generic point is in
$Z^\circ$, and its corresponding divisor on $X$ lies
in $Y$.  The two discrete valuation rings are identical.
Consequently every cover counted by $Q(X,Y)$ restricts
to a cover of $W$ tame at every codimension-one component
of $Z\setminus W$.  We therefore obtain a surjection
\[
                 Q\bigl(Z,E\cup q(Y)\bigr)
                         \longrightarrow Q(X,Y).
\]
The source is finitely generated by the projective case
of Proposition~\ref{D:fixed-boundary-positive}.  This
proves the result.  The codimension-two conclusion follows
because there are then no original boundary inertia groups
to kill in defining $Q(X,Y)$.
\end{proof}

\begin{remark}
The hypothesis is a morphism from $X$ to a projective target.
A projective alteration mapping onto $X$ has the opposite
direction and does not supply this hypothesis.  The proof
uses only unchanged codimension-one DVRs on the target;
it imposes no tameness condition at exceptional divisors
of $X$ whose image has codimension at least two.
\end{remark}
